\section{Evidence Convergence and Handler Routing}
\label{sec:evidence-convergence}

\subsection{Three routing sources and one acceptance observer}

File-type evidence serves two different decisions. Before reconstruction, three
bounded sources select at most one candidate handler. After reconstruction, a
separate complete-parser registry determines whether the candidate belongs to
one observed format profile. Routing is therefore a work-selection decision.
It is never authority to publish.

Let $\mathcal{K}$ be the finite registry of concrete format profiles. Profiles
distinguish materially different grammars, including static PNG from APNG,
static WebP from animated WebP, and constrained WebM from general Matroska. A
versioned relation $\equiv_p$ maps declared aliases such as \texttt{jpg} and
\texttt{jpeg}, or the M4A media types \texttt{audio/mp4},
\texttt{audio/x-m4a}, and \texttt{audio/m4a}, to one profile. Policy
admits $\mathcal{K}_p\subseteq\mathcal{K}$ for positive routing. For an input
$a=(x,n,m,d)$, the routing sources are
\begin{align}
  E_p(n)&\subseteq\mathcal{K}, &
  M_p(m)&\subseteq\mathcal{K}, &
  B_p(x)&\subseteq\mathcal{K}.
  \label{eq:evidence-functions}
\end{align}
$E_p$ interprets the normalized final suffix of the display name. $M_p$
parses the supplied media-type essence and its registered aliases. $B_p$
performs bounded byte recognition. Structural header checks belong to $B_p$
rather than constituting a fourth routing source. Each source returns a set and
a typed status rather than a scalar confidence score.

An evidence report is $\epsilon_i=(S_i,s_i,c_i)$, where
$S_i\subseteq\mathcal{K}$ is the candidate set, $s_i$ is the source status, and
$c_i$ contains bounded counters and reason codes. Status distinguishes an
absent supplied value, unknown alias, malformed syntax, budget exhaustion,
recognition failure, and successful completion. An empty candidate set alone
does not distinguish these cases and is never interpreted as evidence of
safety.

\begin{table}[t]
\centering
\caption{Evidence roles and their authority in the method. Complete parsing is
reserved for source resolution and output acceptance, not counted as a fourth
routing source.}
\label{tab:evidence-strength}
\small
\begin{tabularx}{\textwidth}{L{0.22\textwidth}L{0.22\textwidth}YY}
\toprule
Evidence operation & Stage & Permitted use & Fail-closed interpretation \\
\midrule
Name evidence $E_p$ & Routing & Candidate profile and dangerous-suffix policy & Unknown or conflicting suffix gives no byte authority \\
Supplied type $M_p$ & Routing & Caller intent and provisional compatibility & Missing, malformed, wildcard, or conflicting type is policy-visible \\
Bounded bytes $B_p$ & Routing & Magic and minimum structural-header evidence & Bare, conflicting, or exhausted recognition cannot authorize delivery \\
Complete source parse & Selected reconstruction path & Account for source boundaries and resolve ambiguity required by the chosen method & Failure or unaccounted ranges prevent structural copying \\
Complete output parsers & Acceptance & Whole-object profile observation from offset zero & Prefix success, ambiguity, or parser exhaustion prevents acceptance \\
\bottomrule
\end{tabularx}
\end{table}

No source follows a last-writer-wins rule. A compatible byte probe does not
erase an incompatible supplied type or a dangerous name. HTTP semantics also
distinguish declared type from content inspection~\cite{rfc9110}. The decision
record retains every source, status, disagreement, and resolution step.

\subsection{Normalization and compatibility}

The display name $n$ is not a filesystem path. The normalization function
$N_p(n)$ accepts one basename and rejects path separators, alternate-stream
syntax, dot components, control characters, bidirectional formatting controls,
reserved device names, and trailing platform-sensitive syntax. It records all
suffix components. A forbidden executable, script, archive, shortcut, package,
or macro-bearing nonfinal suffix causes rejection even when the final suffix is
an admitted media extension. Extension comparison uses ASCII case folding and
does not silently remove trailing dots or spaces.

The supplied media type is parsed after syntactic parameter removal. Wildcards
and malformed strings are not positive evidence. A versioned alias table maps
accepted spellings to profiles without treating a complete top-level media
family as supported.

Compatibility is the explicit relation
$\mathsf{Compat}_p\subseteq\mathcal{K}\times\mathcal{K}$. It contains aliases
and narrowly stated container relationships. For example, bounded ISO base
media recognition can identify a generic container before a parent-aware parser
establishes whether the object is audio-only or includes video. This is
provisional compatibility. It is not equivalence among all MP4-family profiles.

\begin{definition}[Input convergence]
\label{def:input-convergence}
Let $I_p(a)=\{\epsilon_E,\epsilon_M,\epsilon_B\}$. Input evidence converges on
$k\in\mathcal{K}_p$ when every present nonempty candidate set contains a profile
compatible with $k$, no report identifies an incompatible or forbidden profile,
every policy-required report is present, and no recognizer terminates by budget
exhaustion. The predicate is written $\mathsf{Conv}_p(a,k)$.
\end{definition}

Strict evaluation treats unresolved disagreement, missing required supplied
type, unknown required name evidence, and indeterminate byte recognition as
non-routable. Diagnostic profiles may record warnings, but warning and dry-run
outcomes are outside the accepted scientific relation.

\subsection{Routing through a singleton eligible set}

Define the eligible route set
\begin{equation}
  \mathcal{C}_p(a)=
  \{k\in\mathcal{K}_p:N_p(n)\land\mathsf{Conv}_p(a,k)\}.
  \label{eq:eligible-routes}
\end{equation}
The routing function is then genuinely partial:
\begin{equation}
  \mathsf{Route}_p(a)=
  \begin{cases}
    k, & \mathcal{C}_p(a)=\{k\},\\
    \bot, & |\mathcal{C}_p(a)|\neq 1.
  \end{cases}
  \label{eq:routing-function}
\end{equation}
The value $\bot$ means that no unique route is authorized. Depending on the
failure boundary, it becomes a blocking verdict or a typed error. Neither case
can expose deliverable bytes.

After selection, a handler may perform a complete source parse or a semantic
decode. Its evidence record must establish
$\mathsf{SourceResolved}_p(a,k,\eta)$. For structural reconstruction, this
predicate requires complete source accounting, explicit treatment of every
copied range, and resolution of ambiguity relevant to the retained bytes. For
semantic reconstruction it requires a bounded successful decode of the selected
interpretation. The predicate is part of acceptance. It is not another source
that competes in routing.

\begin{property}[Routing is non-authorizing]
\label{prop:routing-nonauthorizing}
There are inputs for which $\mathsf{Route}_p(a)=k$ but no deliverable artifact
exists. A recognized header can be followed by truncation, reconstruction can
fail, source resolution can remain incomplete, or a candidate can fail an
output gate. Routing success therefore implies only unique work selection.
\end{property}

\subsection{Policy-independent complete output observations}

Let $\mathcal{O}_v$ be the complete-parser observation registry fixed by
registry version $v$. It includes admitted profiles and forbidden profiles that
remain relevant as negative ambiguity evidence. For candidate bytes $y$, parser
$k$ returns either failure or
\begin{equation}
  \rho_k(y)=(k,o,t,C,q),
  \label{eq:parse-report}
\end{equation}
where $o$ is the start offset, $t$ is the terminal offset, $C$ is the parsed
component tree, and $q$ contains structural facts and resource counters. The
observation set is
\begin{equation}
  \mathsf{Full}_v(y)=
  \{k\in\mathcal{O}_v:\exists C,q\ .\
    \rho_k(y)=(k,0,|y|,C,q)\}.
  \label{eq:full-parser-set}
\end{equation}
Crucially, $\mathsf{Full}_v$ does not remove a parse because a policy forbids a
component. Policy is applied later to the stable parser report. Tightening a
policy therefore cannot hide an alternate complete interpretation.

A successful reconstruction returns the typed candidate
\begin{equation}
  c=(y,n_{\mathrm{out}},m_{\mathrm{out}},k_h,\ell,\eta),
  \label{eq:candidate-artifact}
\end{equation}
where $k_h$ is the handler profile, $\ell$ is the performed level, and $\eta$
is the evidence record. Output convergence for $(c,k)$ requires
$\mathsf{Full}_v(y)=\{k\}$, $k_h$ compatible with $k$, and both the generated
name and declared output type to map to $k$. A successful prefix parse never
enters $\mathsf{Full}_v(y)$.

This rule distinguishes four mechanisms. Prefix camouflage fails the
offset-zero requirement. A trailing object or appended data fails exact
consumption. A true dual-format polyglot produces at least two incompatible
complete observations of the same whole octet string. A nested active object is
handled by the location-aware component policy. Structural reconstruction
cannot claim to neutralize an alternate interpretation inside retained codec
ranges. Semantic reconstruction supports a narrower statement only when a
bounded decoder supplies semantic state, source container and compressed
payload ranges are not copied, and the newly encoded candidate has one complete
observed profile.

\begin{property}[Unique observed output profile]
\label{prop:unique-output-profile}
If output convergence holds for $(c,k)$, no other profile in $\mathcal{O}_v$
has a complete incompatible parse of $c.y$. This is a registry-relative design
invariant. It does not exclude an unmodeled grammar or a common parser defect.
\end{property}

The remaining policy predicates operate on this fixed observation and typed
candidate. They are defined in \cref{sec:formal-method}.
